\usepackage[T1]{fontenc}
\usepackage{mathpazo}
\usepackage[scaled=0.92]{helvet}
% The book's fonts are T1-encoded Palatino and Helvetica, so LuaLaTeX sends a character outside T1 to the slot of
% its code point: × (U+00D7) printed as Œ, and −, ² and Σ were dropped or replaced ("ranges from −1 to +1" read
% "from 1 to +1"). Each is mapped to the text symbol of the companion encoding (TS1), and Σ, which T1 and TS1
% lack, to TeX Gyre Pagella, the OpenType Palatino.
\usepackage{newunicodechar}
\newunicodechar{×}{\texttimes}
\newunicodechar{−}{\textminus}
\newunicodechar{²}{\texttwosuperior}
\ifdefined\newfontfamily
  \newfontfamily\aagreek{TeX Gyre Pagella}
  \newunicodechar{Σ}{{\aagreek Σ}}
\fi
\usepackage{microtype}
\usepackage{booktabs}
\usepackage{longtable}
% The book's tables are longtables (Quarto and Pandoc write every table that way), which can break at any
% row and never move. Four changes:
% - A captioned table no taller than 85% of the text height floats like a figure: it stays where it is if
%   it fits there, otherwise it goes to the top of the next page and the text that follows fills the
%   space (\fps@table below). When a longtable ends, it holds the whole table, its head with the caption
%   in \LT@firsthead, its rows in box 0 and its bottom rule in \LT@lastfoot, and calls \LT@start to set
%   it on the page; \LT@start is wrapped to measure those boxes and, for such a table, to put them in a
%   box of its own instead, which \endlongtable then sets as a float (\aa@LT@float). The caption,
%   number, label and link target are the ones longtable made.
% - A longer table breaks across pages as before, with its header repeated, but never starts in the last
%   quarter of a page, where its caption and a row or two would sit alone (\aa@LT@long).
% - longtable resets LaTeX's room for text on the page (\@colroom) to the full page when a table starts
%   and never sets it back. On a page whose top holds a figure, the next figure after the table then let
%   the text run past the foot of the page (Table 15.2 and Figure 15.2 in the print interior). At the end
%   of a table the room is set back to what the page allows (\vsize: the full page again if it broke).
% - A figure that follows a table which broke across pages could go at the top of the table's last page,
%   between its two parts (Figure 1.4 inside Table 1.3). When a table breaks, no figure goes at the top
%   of the page where it ends: the figure follows the table on that page if it fits, or moves to the next.
\usepackage{etoolbox}
\makeatletter
\newif\ifaa@LTbroke
\newif\ifaa@LTfloat
\newbox\aa@LTbox
\newskip\aa@LTpost
\newdimen\aa@LTheight
\def\aa@LTtable{table}
\let\aa@LT@start@orig\LT@start
\def\LT@start{%
  \aa@LTheight\ht\z@ \advance\aa@LTheight\dp\z@
  \advance\aa@LTheight\ht\ifvoid\LT@firsthead\LT@head\else\LT@firsthead\fi
  \advance\aa@LTheight\dp\ifvoid\LT@firsthead\LT@head\else\LT@firsthead\fi
  \advance\aa@LTheight\ht\ifvoid\LT@lastfoot\LT@foot\else\LT@lastfoot\fi
  \edef\aa@LTcap{\LTcaptype}%
  \let\aa@LTnext\aa@LT@long
  \ifx\aa@LTcap\aa@LTtable
    \ifdim\aa@LTheight>0.85\textheight\else \let\aa@LTnext\aa@LT@float \fi
  \fi
  \aa@LTnext}
% A short table: its parts go into \aa@LTbox; the end penalty that would call longtable's output
% routine becomes an ordinary one, and \endlongtable (below) sets the box as a float.
\def\aa@LT@float{%
  \global\setbox\aa@LTbox\vbox{%
    \MakeLinkTarget{\LTcaptype}%
    \ifvoid\LT@firsthead\copy\LT@head\else\box\LT@firsthead\fi
    \unvbox\z@
    \ifvoid\LT@lastfoot\copy\LT@foot\else\box\LT@lastfoot\fi}%
  \global\setbox\LT@foot\box\voidb@x
  \mathchardef\LT@end@pen\z@
  \global\aa@LTpost\LTpost \global\LTpost\z@skip
  \global\aa@LTfloattrue}
% A long table: not in the last quarter of a page, unless it follows a heading, which must keep it.
\def\aa@LT@long{%
  \if@nobreak\else
    \endgraf\penalty\z@
    \ifdim\pagegoal<\maxdimen
      \dimen@\pagegoal \advance\dimen@-\pagetotal
      \ifdim\dimen@<0.25\textheight \vfil\penalty-\@M \fi
    \fi
  \fi
  \aa@LT@start@orig}
% Every cell sits at the top left of its row. Pandoc sets each header cell of a table with column widths
% in \begin{minipage}[b], which drops a one-line header to the last line of a header that wraps; inside
% a longtable a minipage is set [t] instead, like the body cells.
\let\aa@minipage\minipage
\def\aa@minipage@t[#1]{\aa@minipage[t]}
\AtBeginEnvironment{longtable}{\def\minipage{\@ifnextchar[\aa@minipage@t\aa@minipage}}
\pretocmd{\LT@output}{\ifnum\outputpenalty<-\@Mi\else\global\aa@LTbroketrue\fi}{}
  {\PackageWarning{book-pdf}{Could not patch longtable's output routine}}
\apptocmd{\endlongtable}{%
  \global\@colroom\vsize\ifaa@LTbroke\global\@topnum\z@\fi\global\aa@LTbrokefalse
  \ifaa@LTfloat
    \global\aa@LTfloatfalse \global\LTpost\aa@LTpost
    \begin{table}\box\aa@LTbox\end{table}%
  \fi}{}
  {\PackageWarning{book-pdf}{Could not patch the end of longtable}}
\makeatother
\usepackage{array}
\usepackage{caption}
\usepackage{graphicx}
\usepackage[table]{xcolor}
\usepackage{eso-pic}

% Code looks the same in every block and inline, as in HTML: on the gray of the highlighted SQL
% blocks (shadecolor, in Quarto's Shaded box). Formula blocks (Excel, DAX, M: unhighlighted, so
% verbatim) sit in the same Shaded box; inline code (\aacode, wrapped around \texttt by
% filters/pdf-code.lua) is highlighted with lua-ul, which, unlike a box, still lets a long formula
% break at the end of a line.
\usepackage{framed}
\usepackage{luacolor}
\usepackage{lua-ul}
\providecolor{shadecolor}{RGB}{241,243,245}
\makeatletter
\@ifundefined{Shaded}{\newenvironment{Shaded}{\begin{snugshade}}{\end{snugshade}}}{}
% \leavevmode starts the paragraph first: lua-ul's \highLight in vertical mode (code at the start of a
% table cell, as in the first column of Table 10.2) set the code a line below the other cells of its row.
\DeclareRobustCommand{\aacode}[1]{\leavevmode\highLight[shadecolor]{\hbox{}\nobreak\hskip1.5pt minus 1.5pt\nobreak
  #1\nobreak\hskip1.5pt minus 1.5pt\hbox{}}}
\AtBeginDocument{\@ifundefined{pdfstringdefDisableCommands}{}{\pdfstringdefDisableCommands{\let\aacode\@firstofone}}}
% Code stays in the order of the text, unlike figures and tables. filters/pdf-keep-together.lua asks
% for room (\aaNeedspace) before each code block and its lead-in, so a query of up to 25 lines moves
% whole to the next page instead of splitting, and a longer one leaves at least six lines on its first
% page. Waiting figures and tables are not placed at the top of a page that a code block starts or
% continues on: they would take the room the block was moved there for (Figure 9.9 split the query of
% Tutorial 9.3, Step 2), or sit between the two parts of a long block (Figure 10.8 once sat inside the
% first query of Tutorial 10.3). They go to the next page instead. \aaNeedspace is needspace's
% \Needspace* with that one difference.
\newif\ifaa@incode
\AtBeginEnvironment{Shaded}{\global\aa@incodetrue}
\AfterEndEnvironment{Shaded}{\global\aa@incodefalse}
\let\aa@startcolumn\@startcolumn
\def\@startcolumn{\ifaa@incode\global\@colroom\@colht\else\expandafter\aa@startcolumn\fi}
\newcommand{\aaNeedspace}[1]{%
  \par\penalty-100
  \begingroup
    \setlength\dimen@{#1}%
    \dimen@ii\pagegoal \advance\dimen@ii-\pagetotal
    \ifdim\dimen@>\dimen@ii \global\aa@incodetrue \break \global\aa@incodefalse \fi
  \endgroup}
% A tutorial step's label ("Step 5. Add the cards.") is a paragraph, not a heading, so the page could break right
% after it and leave it alone at the foot. The same filter puts \aaKeepLabel before each label: when the page
% cannot hold the label and the first lines of the step, the label starts the next page. Unlike \aaNeedspace, it
% adds no penalty that would invite a break, only \penalty0, a break the paragraphs allow anyway, so that a page
% already full is sent out before the room is measured (without it, the page after a full one held only the last
% paragraph before the label). Waiting figures and tables stay free to go first.
\newcommand{\aaKeepLabel}[1]{%
  \par\penalty0
  \begingroup
    \setlength\dimen@{#1}%
    \dimen@ii\pagegoal \advance\dimen@ii-\pagetotal
    \ifdim\dimen@>\dimen@ii \break \fi
  \endgroup}
\makeatother

% Long Excel formulas in code blocks wrap after commas and asterisks, with a continuation mark,
% instead of running into the margin, and long lines of highlighted SQL (Highlighting) wrap at
% spaces. Deferred to \begin{document} because a package loaded later in Quarto's template
% redefines the verbatim environment.
\usepackage{fvextra}
% fvextra's continuation mark, \hookrightarrow in math mode, became U+21AA in a font that lacks it (the log read
% "Missing character: There is no ↪", and the PDF printed no mark), so the mark comes from TeX Gyre Pagella.
\ifdefined\aagreek
  \fvset{breaksymbolleft={\tiny{\aagreek ↪}}}
\fi
\AtBeginDocument{%
  \renewenvironment{verbatim}{\VerbatimEnvironment\begin{Shaded}%
    \begin{Verbatim}[breaklines,breakafter={,*},breakaftersymbolpre={}]}{\end{Verbatim}\end{Shaded}}%
  \RecustomVerbatimEnvironment{Highlighting}{Verbatim}{commandchars=\\\{\},breaklines,breakaftersymbolpre={}}}

\definecolor{PrimaryBlue}{HTML}{1A5276}
% A chapter's learning objectives (filters/learning-objectives.lua) sit in a box like the website's
% (.learning-objectives in styles/book.scss): a blue tint, rounded corners, a blue bar on the left. Quarto
% loads tcolorbox for the callouts; \tcbuselibrary works whichever loads it first.
\usepackage{tcolorbox}
\tcbuselibrary{skins,breakable}
% Its background is a named color so that the print interiors (print/book-print-*.tex) can redefine it.
\colorlet{ObjectivesTint}{PrimaryBlue!6}
\newtcolorbox{aaobjectives}{enhanced, breakable, colback=ObjectivesTint, colframe=PrimaryBlue,
  boxrule=0pt, leftrule=3pt, arc=2mm, left=4mm, right=4mm, top=1mm, bottom=2mm,
  before skip=6pt, after skip=10pt}
% A box never splits across two pages: a callout (In Practice, Watch out, Connecting the Dots) or the
% learning objectives that does not fit at the foot of a page moves whole to the next one. Quarto asks
% for breakable callouts, so the breakable style is emptied for every box; the longest callout is about
% a dozen lines, far less than a page.
\tcbset{breakable@true/.code={}}
% With boxes kept whole, a page can end early; leave that space at the foot of the page instead of
% stretching the gaps between paragraphs and headings to fill it.
\raggedbottom
% No paragraph leaves a single line at the foot of a page or carries one over to the top of the next.
\clubpenalty=10000
\widowpenalty=10000
\displaywidowpenalty=10000
\definecolor{Teal}{HTML}{17A589}
\definecolor{Coral}{HTML}{E74C3C}
\definecolor{Amber}{HTML}{F39C12}
\definecolor{Ink}{HTML}{2C3E50}
\definecolor{LightGray}{HTML}{F2F3F4}
\definecolor{BorderGray}{HTML}{D5D8DC}

\color{Ink}
\setkomafont{disposition}{\normalfont\sffamily\bfseries\color{PrimaryBlue}}
\setkomafont{chapter}{\Huge\sffamily\bfseries\color{PrimaryBlue}}
\setkomafont{section}{\Large\sffamily\bfseries\color{PrimaryBlue}}
\setkomafont{subsection}{\large\sffamily\bfseries\color{PrimaryBlue}}
\setkomafont{captionlabel}{\sffamily\bfseries\color{PrimaryBlue}}
\setkomafont{caption}{\small}

% Figure placement: a figure goes where the text calls it if it fits there, otherwise at the top of a
% page, with the text that follows filling the page before it. A figure up to 85% of the text height
% can share a page with text, and a page of figures alone is made only when they fill at least 80% of
% it, so a large figure no longer sits alone on a page while its paragraph runs on before and after it
% (LaTeX's defaults, 70% at the top and float pages half full, put Figure 1.4 on a page by itself).
% Tables that float (see longtable at the top) follow the same rules.
\makeatletter
\renewcommand{\fps@figure}{htp}
\renewcommand{\fps@table}{htp}
\makeatother
\renewcommand{\topfraction}{0.85}
\renewcommand{\textfraction}{0.1}
\renewcommand{\floatpagefraction}{0.8}
\setcounter{topnumber}{2}
\setcounter{totalnumber}{3}
% A section break that places the waiting floats first (the filter's new-page headings, as the groups of the
% quick reference appendix, which is nearly all tables): the floats left over fill pages of their own, and
% they stack from the top of each page instead of floating in its middle (LaTeX's glue on such a page is
% 0pt plus 1fil above and below and 8pt plus 2fil between, which spreads them over it). The glue is set
% inside the group, so every other float page keeps the default.
\makeatletter
\newcommand{\aaClearpageTop}{{%
  \setlength{\@fptop}{0pt}\setlength{\@fpsep}{12pt}\setlength{\@fpbot}{0pt plus 1fil}%
  \clearpage}}
\makeatother

% Running heads (author's request 2026-10-12): on a left-hand page the Part, flush left ("Part II  Accounting
% Analytics with Excel"), and on a right-hand page the chapter, flush right ("Chapter 5  Data Preparation and
% Cleaning in Excel"; "Appendix A  ..." for an appendix; the title alone for an unnumbered chapter, such as
% the preface or a case). The first page of a chapter and a Part's title page keep the page style plain, with
% no head, and the page numbers keep the foot they had: the new page style is plain with a head.
%
% The marks live in two mark classes of their own, not in \leftmark and \rightmark, because Pandoc writes
% \markright after every unnumbered section (Learning Objectives, Opening Scenario, ...), which would put those
% titles in the head. KOMA's \partmark and \chaptermark, called by \part and \chapter, set the marks of numbered
% Parts and chapters, and \markboth, which Pandoc writes after each unnumbered chapter, those of the
% unnumbered ones. A chapter outside the Parts (the preface, the closing thoughts, the appendices) must not
% show the Part before it: Quarto writes \bookmarksetup{startatroot} before such a chapter, and \appendix
% starts the appendices, so each sets a flag that the next chapter's mark uses to clear the Part. A left-hand
% page outside the Parts has no head. The head sits in the top margin, \headsep above the text, so the print
% interior (print/book-print.tex) makes the top margin 0.8in to keep it inside the printer's safe area.
\makeatletter
\NewMarkClass{aapart}
\NewMarkClass{aachap}
\newif\ifaa@outsidepart
\AtBeginDocument{%
  \let\aa@bookmarksetup\bookmarksetup
  \renewcommand*{\bookmarksetup}[1]{\aa@bookmarksetup{#1}\ifstrequal{#1}{startatroot}{\aa@outsideparttrue}{}}}
\apptocmd{\appendix}{\aa@outsideparttrue}{}{}
\renewcommand*{\partmarkformat}{\textbf{\partname\ \thepart}\quad}
\renewcommand*{\chaptermarkformat}{\textbf{\chapapp\ \thechapter}\quad}
\newcommand*{\aa@leavepart}{\ifaa@outsidepart\InsertMark{aapart}{}\aa@outsidepartfalse\fi}
\renewcommand*{\partmark}[1]{\aa@outsidepartfalse\InsertMark{aapart}{\partmarkformat #1}}
\renewcommand*{\chaptermark}[1]{\aa@leavepart\InsertMark{aachap}{\chaptermarkformat #1}}
\let\aa@markboth\markboth
\renewcommand*{\markboth}[2]{\aa@markboth{#1}{#2}\aa@leavepart\InsertMark{aachap}{#2}}
\newcommand*{\aa@headfont}{\sffamily\footnotesize\color{Ink}}
\newcommand*{\ps@aaheadings}{%
  \ps@plain
  \renewcommand*{\@evenhead}{\hb@xt@\textwidth{\aa@headfont\strut\FirstMark{aapart}\hfil}}%
  \renewcommand*{\@oddhead}{\hb@xt@\textwidth{\aa@headfont\hfil\strut\FirstMark{aachap}}}}
\pagestyle{aaheadings}
\geometry{headheight=12pt, headsep=0.12in}
\makeatother

\captionsetup[figure]{position=bottom, justification=raggedright, singlelinecheck=false}
\captionsetup[table]{position=top, justification=raggedright, singlelinecheck=false}
\setlength{\emergencystretch}{3em}
% A higher tolerance lets a line with an unbreakable inline formula be set loosely
% instead of running into the margin.
\tolerance=3000
% A long code block that breaks across pages can leave its page a few points too full (the
% Shaded box's padding: about 9pt for a long DAX block in Chapter 14); accept that silently
% rather than report it as an overfull box.
\vfuzz=10pt
% Part entries in the table of contents set their page numbers in large bold type, and a
% three-digit number is wider than KOMA's default page-number box.
\RedeclareSectionCommand[tocpagenumberwidth=2em]{part}
\renewcommand{\arraystretch}{1.2}

\makeatletter
\renewcommand{\maketitle}{%
  \begin{titlepage}
    \thispagestyle{empty}
    \IfFileExists{visuals/cover/cover.pdf}{%
      \AddToShipoutPictureBG*{%
        \AtPageLowerLeft{%
          \includegraphics[width=\paperwidth,height=\paperheight]{visuals/cover/cover.pdf}%
        }%
      }%
      \null
    }{%
      \centering
      \vspace*{0.28\textheight}
      {\Huge\sffamily\bfseries\color{PrimaryBlue}\@title\par}
      \vspace{1.5rem}
      {\Large\sffamily\@author\par}
      \vfill
    }%
  \end{titlepage}
  \clearpage
}
\makeatother

% The title page (recto, after the cover and a blank page) and the copyright page (its verso), which
% filters/book-pages.lua writes from front-matter/_book-pages.qmd between the cover and the table of contents.
\newenvironment{aatitlepage}{%
  \cleardoubleemptypage
  \thispagestyle{empty}%
  \pdfbookmark[0]{Title Page}{book-title-page}%
  \begin{center}\sffamily\color{Ink}
  \vspace*{0.14\textheight}}{%
  \end{center}
  \clearpage}
\newcommand{\aabooktitle}[1]{{\fontsize{32}{38}\selectfont\bfseries\color{PrimaryBlue}#1\par}\vspace{0.7em}}
\newcommand{\aabooksubtitle}[1]{{\LARGE #1\par}\vspace{1.6em}{\color{Amber}\rule{1.25in}{1.5pt}}\par\vspace{1.6em}}
\newcommand{\aabookauthor}[1]{{\Large #1\par}\vfill}
\newcommand{\aabookedition}[1]{{\large #1\par}\vspace{0.3em}}
\newcommand{\aabookrevision}[1]{{\large #1\par}\vspace{2.5em}}
\newcommand{\aabookpublisher}[1]{{\large\bfseries #1\par}\vspace*{0.04\textheight}}
\newenvironment{aacopyrightpage}{%
  \thispagestyle{empty}%
  \pdfbookmark[0]{Copyright}{book-copyright-page}%
  \vspace*{\fill}
  \noindent\begin{minipage}{0.72\textwidth}\footnotesize\raggedright
  \setlength{\parindent}{0pt}\setlength{\parskip}{0.7\baselineskip}}{%
  \par\end{minipage}
  \clearpage}
